% Chapter 4a discussion --- this chapter's share of draft \S17 per the cut
% list in Project_specification.md \S4.1: assay predictions + single-cell
% open problems + biological grounding (Grant_paper.pdf \S4, translated) +
% forward pointer to Chapter 4b. Fitting consequences, forward connections
% and the population open problems belong to Chapter 4b's discussion.

\section{Discussion}
\label{sec:discussion}

This chapter set out to complete the single-cell theory of the
birth--death--catastrophe process, and that programme is now done: the
single-cell description is complete in distribution --- geometric states at
every time, a geometric quasi-stationary limit, and a burst-size law
identical to it --- with the conditional laws of burst time and burst size,
the multi-founder extension, and the chained-transfer model alongside. What
remains is to say what the results imply for measurement, where they have
already been felt in real systems, and what is still open on the
single-cell side. The population consequences --- fitting, thresholds, and
the comparison of bursting with budding --- are taken up in Chapter~4b.

\subsection{Single-cell assay predictions for lytic systems}
\label{sec:discussion-assay}

For a genuinely lytic system --- \yp{}, \ft{}, lytic phage --- a
limiting-dilution assay (one infected cell per well, counts followed in
time) measures the single-cell process directly, with no population model
interposed. The number of new infectious particles released per infected
cell is an integer-valued random variable, and its full distribution, not
merely its mean, is what determines whether a small inoculum grows into a
detectable infection \cite{williams2024reproduction}; assays of this design
are where that distribution becomes observable. The theory makes five
predictions, all of them parameter-free given $(\lambda,\mu,\delta)$:
\begin{enumerate}
\item a fraction $b$ of wells show no release ever --- the mass of internal
extinction;
\item the cumulative release in a well is a step function: zero, then a
single terminal jump;
\item the jump times follow the density $\delta J(t)/(1-b)$;
\item the jump sizes follow the geometric law with ratio $1/a$ --- the
quasi-stationary distribution, by Corollary~\ref{cor:burst-qsd};
\item the between-well variance of total release is $\mathrm{Var}(W_\infty)$
of \eqref{eq:EW2}, a quantity no deterministic constant-rate model can
produce at all.
\end{enumerate}
One honest caveat: these predictions are too strong for a budding virus as
literally stated --- intermittent budding is not a single terminal burst ---
and they should be used for lytic hosts, or as the bursting endpoint of the
spectrum of release models that Chapter~4b lays out.

\subsection{The same theory at two extremes}
\label{sec:two-extremes}

The single-cell theory of this chapter has already been fitted to two
intracellular pathogens whose parameter estimates sit at opposite ends of
what the model allows, and the contrast is worth recording here because it
shows how much of the behaviour is fixed by $(\lambda,\mu,\delta)$ alone.

For \ft{} strain SCHU S4, modelling of murine challenge data yielded the
estimates $\lambda=0.15\,\mathrm{h}^{-1}$,
$\mu=0.01\,\mathrm{h}^{-1}$ and
$\delta=1.5\times10^{-4}\,\mathrm{h}^{-1}$
\cite{carruthers2020stochastic}. In the model with these parameters, fewer
than $7\%$ of infected cells fail to burst --- the internal-extinction
probability $b$ is small --- and the burst sizes of the cells that do
rupture are geometrically distributed with typical sizes in the hundreds of
bacteria. That prediction sits comfortably with the observed ability of
doses as low as 10 colony-forming units to establish infection
\cite{oyston2004tularaemia}: each successful rupture floods the local
environment with enough bacteria that a small inoculum needs only a few
rounds of the cycle to reach large numbers. Recent video recordings of
\ft{}-infected macrophages, showing en-masse release events and direct
counts of near-rupture payloads, provide observational support for the
burst-size prediction.

For \textit{Bacillus anthracis}, approximate Bayesian computation against
published experimental time-series produced the very different estimates
$\lambda\approx0.64\,\mathrm{h}^{-1}$,
$\mu\approx1.64\,\mathrm{h}^{-1}$ and
$\delta\approx0.04\,\mathrm{h}^{-1}$ \cite{williams2021anthrax}. Here
death dominates birth: the model predicts that more than $96\%$ of
phagocytes clear their intracellular infection without rupturing, and the
cells that do rupture release on average only about $1.6$ bacteria each ---
the value of $\mean{\Kb\mid\text{burst}}=a/(a-1)$ at these parameters. The
consistent picture is that an infectious dose of $8\times10^3$ to
$5\times10^4$ spores is required \cite{williams2021anthrax}, orders of
magnitude above the tularemia figure, and one consistent with a pathogen
whose intracellular dynamics destroy far more than they release.

These are model-based predictions, and the biological claims rest on the
parameter estimates they condition on; but the two cases together show that
the same three-rate theory spans the space from near-certain, large bursts
to near-certain clearance with minimal release, and that the quantities
fitted in each case --- the no-burst fraction $b$, the burst-size law, the
mean conditional burst --- are exactly the assay observables listed above.

\subsection{Open problems}
\label{sec:open-problems}

Three questions remain open on the single-cell side.

\begin{enumerate}
\item \textbf{A conceptual proof that burst size equals QSD.} The algebra
of \S\ref{sec:burst=qsd} is complete; a structural argument --- presumably
size-biasing composed with the burst intensity, or a Yaglom-type limit
\cite{yaglom1947certain} --- would turn a coincidence into a theorem one
can quote.
\item \textbf{Chained joint laws for general $\mu$.} The
$\mu=0$ chain of \S\ref{sec:chained} rests on the event-type decomposition,
which fails once internal extinction is possible. The joint laws of
$(t_k,r_k)$ for $\mu>0$ --- and the joint law of $(t_M,r_M)$ for the full
chain --- remain open; a coefficient machinery with ${}_2F_1$
antiderivatives is the route sketched in \S\ref{sec:chained}.
\item \textbf{Logistic intracellular growth.} A referee will object that
linear birth lets the conditional load grow without bound. The defence is
that the burst hazard is proportional to load and so regulates it --- the
QSD result is precisely the statement that the conditional load stabilises
at $a/(a-1)$ --- but the sensitivity of the burst statistics, and of the
kernels that Chapter~4b will use, to a logistic birth rate remains to be
checked numerically.
\end{enumerate}

\subsection*{What comes next}

Everything in this chapter is single-cell. The released particles have so
far been counted and then forgotten, and that is the approximation the next
chapter removes. In Chapter~4b, the survival and release statistics derived
here become the age-structured kernels of a population-level renewal model:
cells infected at different times contribute to the infected-cell and
free-particle pools according to $\Ifix(a)$ and $g(a)=\delta K(a)$, the
constant-rate models of viral dynamics appear as a projection of that
renewal system, and the bursting strategy is compared against budding at the
level where the comparison matters --- establishment, invasion speed, and
what data can and cannot identify.
